\subsection{Yêu cầu phi chức năng}

Các yêu cầu phi chức năng được thiết lập dựa trên các tiêu chuẩn chất lượng phần mềm nhằm đảm bảo hệ thống không chỉ thực hiện đúng chức năng mà còn vận hành tối ưu, an toàn và dễ dàng mở rộng.

\subsubsection{Hiệu năng và khả năng sử dụng}

\begin{itemize}
    \item \textbf{[NFR1] Hiệu năng phản hồi (Performance Efficiency)} \hfill (\textit{Ưu tiên: Cao})\\
    Thời gian phản hồi của hệ thống cho các tác vụ nghiệp vụ thông thường không vượt quá 3 giây trong điều kiện tải bình thường.

    \item \textbf{[NFR2] Hiệu năng AI Agent (Performance Efficiency)} \hfill (\textit{Ưu tiên: Cao})\\
    Thời gian phản hồi của trợ lý ảo khi thực hiện suy luận, trả lời câu hỏi hoặc gợi ý sản phẩm không vượt quá 5 giây.

    \item \textbf{[NFR3] Tính tinh gọn giao diện (Usability)} \hfill (\textit{Ưu tiên: Cao})\\
    Giao diện phải được thiết kế thân thiện, tối ưu hóa quy trình trải nghiệm để khách hàng có thể hoàn thành đặt hàng cơ bản trong tối đa 5 bước.

    \item \textbf{[NFR4] Khả năng dẫn dắt (Learnability)} \hfill (\textit{Ưu tiên: Trung bình})\\
    Trợ lý ảo có khả năng hướng dẫn và giúp khách hàng mới sử dụng thành thạo các tính năng hệ thống trong tối đa 3 lần tương tác đầu tiên.

    \item \textbf{[NFR5] Tính dễ học (Learnability)} \hfill (\textit{Ưu tiên: Cao})\\
    Người dùng có thể tự làm quen và sử dụng thành thạo các chức năng tìm kiếm, so sánh và thanh toán chỉ sau 2 lần thao tác thực tế.

    \item \textbf{[NFR6] Chất lượng tương tác (Usability)} \hfill (\textit{Ưu tiên: Cao})\\
    Trợ lý ảo sử dụng ngôn ngữ tự nhiên, không lạm dụng thuật ngữ kỹ thuật; đảm bảo ít nhất 80\% phản hồi đạt yêu cầu theo tiêu chuẩn trải nghiệm người dùng.
\end{itemize}

\subsubsection{Độ tin cậy, bảo mật và khả năng bảo trì}

\begin{itemize}
    \item \textbf{[NFR7] Độ ổn định dịch vụ (Reliability)} \hfill (\textit{Ưu tiên: Trung bình})\\
    Hệ thống duy trì hoạt động ổn định với tỷ lệ gián đoạn dịch vụ (downtime) không vượt quá 1\% trong suốt giai đoạn thử nghiệm và vận hành thử.

    \item \textbf{[NFR8] Bảo mật thông tin (Security)} \hfill (\textit{Ưu tiên: Cao})\\
    Bảo vệ thông tin cá nhân bằng cơ chế mã hóa dữ liệu nhạy cảm và tuyệt đối không lưu trữ mật khẩu dưới dạng văn bản rõ ràng (plain text).

    \item \textbf{[NFR9] Khả năng mở rộng (Maintainability)} \hfill (\textit{Ưu tiên: Trung bình})\\
    Kiến trúc hệ thống linh hoạt, cho phép tích hợp thêm các module thanh toán hoặc vận chuyển mới mà không cần thay đổi cấu trúc nền tảng.

    \item \textbf{[NFR10] Hiệu suất quản trị (Performance Efficiency)} \hfill (\textit{Ưu tiên: Thấp})\\
    Giao diện quản trị đảm bảo tốc độ truy xuất và hiển thị các báo cáo tổng hợp dữ liệu lớn trong vòng tối đa 5 giây.

    \item \textbf{[NFR11] Tính toàn vẹn dữ liệu (Functional Consistency)} \hfill (\textit{Ưu tiên: Cao})\\
    Đảm bảo cơ chế đồng bộ hóa dữ liệu tuyệt đối giữa luồng ghi và luồng đọc (CQRS), ngăn chặn tình trạng sai lệch tồn kho hoặc dữ liệu không nhất quán giữa cache và database.
\end{itemize}